\chapter{Discussion}
\label{ch:discussion}

\draftnote{this chapter is a placeholder skeleton built around the
draft-note prompts left in \cref{ch:hyp1,ch:hyp2,ch:hyp3}. Resolve those
prompts first (they ask for specific numbers read off specific figures),
then rewrite this chapter as connected prose rather than a list.}

This chapter returns to the unifying hypothesis of \cref{ch:introduction}
--- that time-lapse GPR can decouple and accurately track subwavelength
material substitutions and mechanical fluid-front movements using
noise-robust multi-dimensional phase-plane regression and sign-bit
time-reversal, even under heavy-tailed Laplace noise --- and integrates the
results of Hypotheses 1--3 to evaluate it directly.

\section{Hypothesis 1: Translation Tracking}

\Cref{sec:hyp1-lateral,sec:hyp1-vertical,sec:hyp1-diagonal} established an
amplitude-based detectability floor in all three translation directions,
and \cref{sec:hyp1-phaseplane} showed the phase-plane estimator remaining
accurate at lateral and vertical scales where that amplitude floor has
already been reached. \draftnote{state the improvement factor explicitly:
floor scale found in \cref{sec:tl-detectability,sec:vtl-detectability,sec:dtl-detectability}
versus the smallest scale at which \cref{sec:tlp-horizontal-validation,sec:tlp-vertical}
still recover $(\Dz,\Dx)$ within tolerance. Note that the diagonal
phase-plane fit itself is still outstanding (\cref{sec:tlp-diagonal}), so
Hypothesis 1's diagonal claim is currently supported only at the amplitude
level --- flag this explicitly rather than overstating the result.}

\subsection{Lateral/Vertical Asymmetry}

\Cref{sec:th-duality} predicted, from first principles, that a migrated
image separates lateral spatial frequency like a prism but not vertical
frequency. \draftnote{state whether this asymmetry was visible only in the
instantaneous-phase analysis (\cref{fig:tlp-instphase-horiz-summary} vs.\
\cref{fig:tlp-instphase-vert-summary}) or also at the amplitude level
(\cref{fig:tl-summary} vs.\ \cref{fig:vtl-summary}), and what that implies
for a field deployment that cares more about one direction than the other.
Discuss whether the diagonal detectability floor of
\cref{sec:dtl-detectability} sits closer to the (tighter) lateral floor or
the (looser) vertical one, as a further test of the same asymmetry.}

\subsection{Hypothesis 1.5: Local Phase-Gradient Methods}

\draftnote{state whether the local, trace-based methods of
\cref{sec:hyp1-h15} (instantaneous phase, spectral-line fitting, cross-phase
spectrograms, localised STFT) gave results consistent with the global WLS
fit wherever both were computed, and whether Hypothesis 1.5 can be
considered supported, exploratory-but-promising, or unresolved given that
no direct quantitative comparison between the two has yet been run.}

\section{Hypothesis 2: Fluid-Front Tracking and Material Change}

\Cref{sec:ff-phaseplane} is the key test of whether the geometry/material-change
decoupling of \cref{sec:th-material-change} survives contact with a
non-rigid, spatially distributed target. \draftnote{state explicitly, for
each migration method, whether the fitted slopes $(\Dz,\Dx)$ stayed near
zero while the intercept $c$ tracked the front advance as predicted --- this
is the central claim of Hypothesis 2 and should be stated as a number, not
left implicit in a figure reference.}

\section{Hypothesis 3: Noise Robustness}

\Cref{sec:hyp3-purenoise} showed that Kirchhoff and Gazdag respond
differently to pure noise (false-coherent bands versus incoherent texture),
and \cref{sec:hyp3-signbit} introduced sign-bit time-reversal as the
noise-robust excitation scheme for back-propagation. \draftnote{state
whether the amplitude-level noisy results of
\cref{sec:hyp3-lateral,sec:hyp3-vertical,sec:hyp3-diagonal,sec:hyp3-fluidflow}
support the claim that ``the right migration technique'' (per Hypothesis 3)
makes noisy detection possible, and be explicit that --- as flagged
repeatedly in \cref{ch:hyp3} --- the quantitative phase-plane-regression
evidence for this currently exists only for the lateral case
(\cref{fig:tlp-noise}); the vertical, diagonal, and fluid-front cases are
amplitude-level only until the outstanding work of \cref{sec:hyp3-gaps} is
completed. Hypothesis 3 should therefore be reported as \emph{partially}
supported, not fully supported, at the current state of the thesis.}

\section{Hypothesis 4: Future Work}
\label{sec:disc-hyp4}

Hypothesis 4 --- that the method generalises to complex synthetic scenes
with multiple, independently-moving scatterers and non-uniform fluid
fronts, and to real field data --- has not been tested in this thesis.
Validating it would require, at minimum: (1) a synthetic model containing
several scatterers at different depths moving in different directions and
amounts simultaneously, to test whether the phase-plane fit (which assumes
a single dominant displacement within its ROI) can be applied locally
enough to separate multiple independent events; (2) a fluid-front geometry
that is not a straight line, to test whether the front-tracking approach of
\cref{ch:hyp2} generalises beyond the idealised straight front used there;
and (3) a real zero-offset field survey, ideally one with a fluid front
advancing away from a borehole under at least partially known conditions,
to test the field pre-processing assumptions (dewow, time-zero correction,
trace re-binning) that the synthetic data in this thesis sidesteps by
construction. None of this exists yet in this project; Hypothesis 4 is
recorded here as the clearest direction for future work, not as a result.

\section{Limitations}

\draftnote{list concrete limitations: synthetic gprMax data only (no field
validation, pending Hypothesis 4); back-propagation migration not run on
the noisy lateral dataset due to computational cost
(\cref{sec:hyp3-lateral}); the domain-size and STFT-scale-labelling
uncertainties flagged in \cref{ch:methodology,sec:tlp-stft}; the fluid-front
model uses a single idealised thin-layer geometry rather than a swept range
of fracture thicknesses or contrasts; the diagonal experiment
(\cref{sec:hyp1-diagonal}) sweeps only a single fixed $2{:}1$ lateral-to-vertical
ratio rather than a range of diagonal angles.}

\section{Outlook for Field Application}

\draftnote{expand with concrete next steps, building on
\cref{sec:disc-hyp4}'s Hypothesis 4 discussion: (1) validate the pipeline on
a real zero-offset field survey with a known, controlled sub-wavelength
displacement (e.g.\ a target on a calibrated micrometre stage) to test the
field pre-processing assumptions of \cref{sec:meth-phaseplane} that the
synthetic data in this thesis sidesteps by construction; (2) complete the
outstanding quantitative noise-robustness work of \cref{sec:hyp3-gaps}; (3)
extend the fluid-front model of \cref{ch:hyp2} to a swept range of fracture
thickness and fluid contrast, and to a front geometry that is not perfectly
straight; (4) quantify the velocity-recalibration procedure needed when
soil or ice moisture genuinely changes between baseline and monitor
surveys, since \cref{sec:meth-phaseplane} noted that this is otherwise
indistinguishable from a true vertical shift $\Dz$.}
